\section{Discussion}
\label{sec:discussion}

\subsection{What is established}
The paper offers three things.
First, a vocabulary for treating mathematical constructions as closures (Sections~\ref{sec:dictionary} and~\ref{sec:packaging-engine}), with the two contracts that make the vocabulary precise: packaging needs a genuine pseudometric, and an operator becomes a packaged law only if it respects the packaging.
Second, machine-checked anchors that pin the calculus exhibit to exact statements: the finite-difference product rule, the limit bridge to the Leibniz rule, the classification of polynomial derivations, and the packaging, descent, and linear route bounds.
Third, numerical diagnostics that were able to fail, together with the controls that make their outcomes interpretable: a matched population for the stencil experiment, null charts and truth checks for the route comparisons, a proof for the prime control regime, and a classical theorem for the endpoint of the positivity toy.

\subsection{Limits of the claims}
\paragraph{Diagnostics are not theorems.}
The stencil experiment shows that, in a fixed population of $\StencilMatchedTotal$ normalized stencils, the Leibniz gate keeps exactly the first-derivative ones, and it explains why through the limiting defects.
It is not a classification of all local operators, and the false-positive search is sampled, not exhaustive.
Fitted exponents, such as $p\approx\HolonomyExponent$ for the coordinate change, summarize a few grids and are not asymptotic orders.

\paragraph{Prime closures are a baseline, not a continuation.}
In the region $\Re(s)>1$ the diagnostics confirm what Proposition~\ref{prop:prime-control} proves.
In the strip they document that naive staging fails, for reasons that are already classical, and they say nothing about the zeros of $\zeta$.

\paragraph{The positivity toy is a motif.}
Its endpoint is the Lee--Yang theorem; its intermediate trend is one sampled system.
The sign constraint it uses is not a general passivity condition.

\paragraph{The dictionary is adopted.}
The six roles organize constructions that are already present in mathematical practice.
They do not supply the comparison spaces, Cauchy conditions, or descent estimates; those must be proved in each case, as Sections~\ref{sec:packaging-engine} and~\ref{sec:exhibit-calculus} do for the calculus exhibit.

\subsection{Next steps}
\paragraph{Analytic closure.}
The calculus exhibit could replace its finite gates with classical consistency and stability hypotheses, and state the regularity classes under which each ledger term vanishes.
The coordinate-change exhibit could connect its observed decay to truncation-order estimates for the interpolation and differencing steps.

\paragraph{Prime ledgers that are valid in the strip.}
A natural next protocol uses the approximate functional equation as an additive staging that is valid in the critical strip, compared against smoothed prime-sum or logarithmic-derivative ledgers in regimes where they are controlled \cite{Titchmarsh1986}.
The present exhibit is the baseline against which such a scheme would be measured.

\paragraph{Formalization.}
The quotient-valued derivation of Proposition~\ref{prop:quotient-derivation} and the convergence argument of Proposition~\ref{prop:prime-control} are short and standard, and both are natural candidates for mechanization alongside the existing bridges.

\subsection{Why the closure lens helps}
Even at this modest scale, the closure lens separates two questions that are easy to conflate: whether a formula for a limit object exists, and whether the staged protocol behind it is stable under refinement and compatible with the operations one wants to perform on it.
The same ledger, packaging, and route-mismatch ideas apply to measure and integration, to limits in functional analysis, to the patching of local descriptions in geometry, and to spectral constructions.

\subsection{Closing}
The slogan ``to count a stone with six birds'' is meant literally.
Counting is a staged discrete protocol; calculus is a stable compression of that protocol under an accounting ledger; and higher layers arise by repeating the packaging with new constraints and new mismatch diagnostics.
The aim of the paper is to make those roles explicit, auditable, and reusable.
